\begin{activity}{La légende du jeu d'échec} \label{echec_et_maths}

\begin{objective}
Analyser un récit historique et faire preuve d'esprit critique.
\end{objective}


La légende la plus célèbre sur l'origine du jeu d'échecs raconte %
l'histoire du roi Belkib, roi des Indes, {3\,000} ans avant notre ère qui %
cherchait à tout prix à tromper son ennui.
 Il promit donc une récompense exceptionnelle à qui lui proposerait une distraction qui le satisferait. 
 
 Lorsque le sage Sissa, fils du Brahmine Dahir, lui présenta le jeu d'échecs, le souverain, enthousiaste, demanda à Sissa ce que celui-ci souhaitait en échange de ce cadeau extraordinaire. 
 Humblement, Sissa demanda au prince de déposer un grain de riz sur la première case, deux sur la deuxième, quatre sur la troisième, et ainsi de suite pour remplir l'échiquier en doublant la quantité de riz à chaque case. 
 Le prince accorda immédiatement cette récompense en apparence modeste, mais son conseiller lui expliqua qu'il venait de signer la mort du royaume car les récoltes de l'année ne suffiraient à s'acquitter du prix du jeu.

\begin{center}
\begin{tikzpicture}
\draw (1,1) rectangle (9,9);
\foreach \x in {1,3,...,8}
\foreach \y in {1,3,...,8}
{ 
\draw[fill=black] (\x,\y) rectangle (\x+1,\y+1);  
\draw[fill=black] (\x+1,\y+1) rectangle (\x+2,\y+2);
}
\end{tikzpicture}
\end{center}

\end{activity}
